% Source: https://artofproblemsolving.com/community/c3h2427111_several_problems_that_i_made
% Stable identity: aops-topic-2427111-q25
Translate so that $O=(0,0)$ and $B=(5,12)$ (this changes neither distances nor slopes); the radius is $OB=13$.

\textbf{Where is $C$?} $C=(c,0)$ for some $c$. The right angle of triangle $OBC$ cannot be at $O$, because $OB$ is not vertical. If it were at $B$, then $(C-B)\cdot(B-O)=0$ gives $5(c-5)-144=0$, $c=\dfrac{169}{5}>13$, so $C$ would be outside the circle. Hence the right angle is at $C$, and $C=(5,0)$, which is inside the circle.

\textbf{The lines.} The tangent $l$ is perpendicular to the radius $OB$, whose slope is $\dfrac{12}5$, so $y=-\dfrac{5}{12}$. The distance between $l$ and the parallel line $t$ is the length of the projection of $\overrightarrow{CB}=(0,12)$ onto the direction of $OB$, which is perpendicular to both lines:
\[x=\dfrac{(0,12)\cdot(5,12)}{13}=\dfrac{144}{13}.\]

\begin{asy}
usepackage("amsmath");
size(5.5cm);
defaultpen(fontsize(10pt));
pair O=(0,0),B=(5,12),C=(5,0);
draw(circle(O,13),gray);
pair d=unit((12,-5));
draw(B-9*d--B+9*d,blue); draw(C-9*d--C+9*d,red);
draw(O--B,dashed); draw(O--C); draw(B--C);
draw(C+(-0.9,0)--C+(-0.9,0.9)--C+(0,0.9));
pair H=C+dot(B-C,unit(B))*unit(B);
draw(C--H,deepgreen+linewidth(1)); draw(H+0.9*unit(-B)--H+0.9*unit(-B)+0.9*d--H+0.9*d);
dot(O); dot(B); dot(C);
label("$O$",O,SW); label("$B$",B,NE); label("$C$",C,S); label("$l$",B+9*d,E,blue); label("$t$",C+9*d,E,red);
label("$x$",(C+H)/2,SW,deepgreen);
\end{asy}

So
\[xy=\dfrac{144}{13}\cdot\left(-\dfrac5{12}\right)=\boxed{-\dfrac{60}{13}}.\]
